\scanpage{6}
\mspage{4}

\textcircled{3} \uline{Descente de} \deleted{\unsure{ses}} \uline{faisceaux}, \deleted{\unsure{d'immersions}} \deleted{\unsure{de morphismes}} \uline{qu. coh.}, \deleted{\uline{et de morphismes}} \inserted{de \deleted{\unsure{morphismes}} sous-schémas}, de morphismes de schémas, de schémas affines.

\deleted{\uline{Prop} $f:S'\mathrel{-}S$ fid plat et qu. cpt, $F$ qu. coh. sur $S$. Alors}

\[
\deleted{\operatorname{Ssf}(F)\mathrel{-}\operatorname{Ssf}(F')\rightrightarrows\operatorname{Ssf}(S'')\quad\text{est exact.}}
\]

\deleted{On est ramené au cas $S$ affine, puis $S'$ affine, et au résultat \unsure{connu|anc.|conn.} d'algèbre commutative.}

\diagram{\begin{tikzcd}[ampersand replacement=\&,row sep=small]
N'\cap M=N_0 \& N' \& {}\\
 M \& M'\arrow[u] \& M''\\
 A\arrow[r,dash] \& A'\arrow[r,dash]\arrow[r,shift right=.5ex,dash] \& A''
\end{tikzcd}}

\uline{Th} Soit $f:S'\mathrel{-}S$ fid plat qu. cpt. Alors $f$ est de descente effective (univ.) pour la catégorie fibrée des Modules qu. coh.

Expliquer ce que cela signifie.

\uline{Cor} \deleted{$f$ est de \unsure{d.}} Soit $F$ qu. coh sur $S$. Alors suite exacte

\[
\operatorname{Ssf}(F)\mathrel{-}\operatorname{Ssf}(F')=\operatorname{Ssf}(F'')
\]

\uline{Corollaire} On a des suites exactes

\deleted{\unsure{``ssf$(S)$''|$\operatorname{Ssf}(S)$}}

\[
\begin{aligned}
 \operatorname{Ss}\deleted{\unsure{F}}f(S)&\mathrel{-}\operatorname{Ssf}(S')\rightrightarrows\operatorname{Ssf}(S'')\\
 \operatorname{Ss\,ou}(S)&\longrightarrow\operatorname{Ss\,o}(S')\rightrightarrows\operatorname{Ss\,o}(S'')\\
 \operatorname{Ss\,rc}(S)&\mathrel{-}\operatorname{Ss\,rc}(S')\rightrightarrows\operatorname{Ss\,rc}(S'')
\end{aligned}
\]

\margin{Peut-on descendre les sous-schémas \unsure{par imm. rétroc.|pas nécess. rétroc.} \unsure{ou ouvertes|ni ouvertes|non ouverts} ?}

\uline{Corollaire} $f$ est de descente (\unsure{univ|univ.}) pour la \unsure{cat-cart|catégorie} fibrée des schémas sur d'autres :

\[
 \operatorname{Hom}_S(X,Y)\mathrel{-}\operatorname{Hom}_S(X',Y')=\operatorname{Hom}_S(X'',Y'')
\]

\diagram{\begin{tikzcd}[ampersand replacement=\&]
 X\times_S Y=Z\arrow[d,dash] \& Z'\arrow[d,dash] \& Z''\arrow[d,dash]\\
 Y\arrow[u,bend left] \& Y' \& Y''\\
 S\arrow[r,dash] \& S'\arrow[r,shift left=.65ex,dash]\arrow[r,shift right=.65ex,dash] \& S''
\end{tikzcd}}

On le sait par le corollaire précédent si $Y$ est qu. sép sur $S$. Mais on se ramène aisément au cas $S,Y$ affines\ldots

On retrouve que $g'$ un iso $\Rightarrow g$ un iso

\uline{Corollaire} Pour que $f'$ une imm. fermée (ouverte, \inserted{arbitraire} \deleted{\unsure{propre}}), il f et \uline{s que} $f$ le soit

\deleted{\unsure{Ces prop gén si $g'$ une imm fermée et $g$} \unsure{aussi ? O.K. si $g$ un morphisme}}

\uline{Corollaire} \deleted{des} d'Algèbres qu. coh.

$\Downarrow$ \deleted{Cor} \uline{Th} \inserted{Si $f$ fid plat qu. cpt, alors} \deleted{\unsure{$f$ est fid}} un morphisme de descente eff. (univ) pour les schémas relatifs \uline{affines}

\uline{Cor} $g$ \deleted{une immersion} affine $\Longleftrightarrow g'$ affine

\[
\begin{aligned}g\ \text{quasi-affine}&\Longleftrightarrow g'\ \text{quasi-affine}\\
 g\ \text{entier (fini, \inserted{fini loc. libre})}&\Longleftrightarrow g'\ \text{entier (fini, fini loc. libre)}
\end{aligned}
\]

\margin{[car $g$ qu. aff $\Longleftrightarrow g$ \unsure{qch|qu. cpt|coh.} et $X\mathrel{-}\operatorname{Spec}g_*(O_X)$ une imm}

% archive sheet number: 5
